\chapter{The Two-Phase Law}\label{ch:twophase}

\NewV\ \emph{This chapter replaces the role a single global objective would
have played. There is no such objective, and the reason is structural.}

\begin{workplan}{\NewV}
\wpgoal{Establish that \textsc{tatiana} alternates between two phases with
different objectives, that this is forced rather than chosen, and that cost
lives in exactly one of them.}

\wpsource{\texttt{TATIANA\_V1\_IDENTITY\_MAP.md}~\S4 (the table), \S4.1 (wake,
the Hodge routing law), \S4.2 (sleep), \S6 (the maintenance cost);
\texttt{MATH\_TOOLBOX.md}~I.5 (predictive coding, $\dot x = -Lx$ is sheaf
diffusion), II.2 ($\gamma(\nu)$, $\gamma_0 \ll 1$).}

\wpsteps{
  \item \textbf{Open with the table}, since it is the whole chapter in one
    view: what drives each phase ($q_t \neq 0$ vs $q_t = 0$), what moves
    ($s$ fast, $\Fs$ slow, $\CC$ grows / $\Fs$ renormalizes, $\CC$ prunes),
    the objective (discharge discord / pay down maintenance cost), and what
    happens to $\mathrm{K}(\mathbb{K})$ (conserved exactly / may decrease).
  \item \textbf{Argue both phases are forced.} The model must be able to run
    with empty input; that requirement alone establishes an offline phase
    exists. Name it rather than deriving it from biology, and only then note
    the neuroscience that agrees.
  \item \textbf{Wake --- the Hodge routing law.} On the $2$-complex,
    $\Cone = \im\cob^0 \oplus \Hone \oplus \im(\cob^1)^*$, orthogonally. State
    and prove the routing claim: \emph{each summand is dischargeable by
    exactly one layer.} $\im\cob^0$ yields to changing $s$ (the fast flow
    $\dot s = -\Lap_{\Fs} s + \text{drive}$, provably, since $\Lap_{\Fs}$ is
    PSD); $\Hone$ does not --- it is holonomy of $\Fs$ around a cycle, and only
    molding or growth touches it; $\im(\cob^1)^*$ flags a bad $2$-cell.
    Include the Lyapunov argument in full.
  \item \textbf{The handoff is a stall, not a threshold.} This is the
    chapter's sharpest point and must not be blurred. When the constrained
    descent has converged and discord remains, it is \emph{proven} that no
    further motion at that level helps. That is the growth trigger, derived
    from the geometry. Contrast explicitly with a rule of the form
    $\rho > c$, and say why the latter would be a hand-set number.
  \item \textbf{Name the one free parameter.} The numerical tolerance on
    ``converged'' is an \Proposed\ engineering choice. State it in the open.
    Do not hide it inside a $\lambda$.
  \item \textbf{Sleep --- renormalize, then down-select.} Apply $\Pi_{O(d)}$
    globally: scale drops, relative structure survives
    (Chapter~\ref{ch:identity}, level 1). Then consult the maintenance cost
    and prune structure that does not pay for itself.
  \item \textbf{State where energy lives and why it lives nowhere else.} One
    term, one phase. Growth adds during wake; pruning removes during sleep;
    they never bid against each other inside a single tick. Give this its own
    paragraph --- it is the entire answer to the objection that the design
    needs a multi-term objective with weights.
  \item \textbf{The maintenance cost.} $\mathcal{U}(\mathbb{K}) =
    L(\text{structure}) + L(\text{observations} \mid \text{structure})$. The
    description-length form is already derived (Chapter~\ref{ch:mdl}); what is
    new here is only its \emph{placement}. Then name the one genuine free
    parameter of the entire design: the exchange rate between the cost of a
    deformation and the cost of a new cell, since MDL needs both in the same
    units. State it in the open and record that if it acquires siblings, the
    design has regressed.
}

\wpdone{A reader can say what moves in each phase, prove that the fast flow
cannot discharge a harmonic component, explain why the growth trigger is not
a threshold, and name the single exchange rate together with the reason a
second one would be evidence of failure.}
\end{workplan}

\begin{workplan}{\Retire}
\wpgoal{\S: Record the refused alternative, so it is not proposed again.}

\wpsource{The predictive-coding discussion in
\texttt{DOCS/Phase I --- Understand the Geometry o.txt} (second half);
refusal recorded in \texttt{TATIANA\_V1\_IDENTITY\_MAP.md}~\S8 DON'T~1 and
\S13.}

\wpsteps{
  \item State the proposal as it was made and at its strongest: a single
    cognitive energy functional $F = E_{\text{pred}} + \lambda_C E_C +
    \lambda_{\Fs} E_{\Fs} + \lambda_s E_s + \lambda_I E_I + \lambda_K E_K$,
    with dynamics choosing whichever admissible transformation most decreases
    it. It is a coherent design and its appeal should be visible.
  \item Give the refusal and its reason: six weights are six hand-set numbers
    trading against each other every tick, which is exactly the failure mode
    the discipline rule against hand-set constants exists to prevent. The
    two-phase separation achieves the same coordination with one exchange
    rate.
  \item Record the standing consequence: \emph{do not add a $\lambda$.} A
    second exchange rate means the two-phase separation has failed and should
    be re-derived, not patched.
  \item Note honestly what the refused proposal got right, since parts of it
    survive under other names: prediction error precision-weighted is $\omega$
    (Chapter~\ref{ch:rho}); the compression term is the MDL law
    (Chapter~\ref{ch:mdl}); the stability term is $\gamma_0 \ll 1$'s
    anti-interference property. What is rejected is the \emph{blending}, not
    the ingredients.
}

\wpdone{Someone rereading the predictive-coding notes can see why they were
not adopted wholesale, and which of their parts already live elsewhere in the
book.}
\end{workplan}
